\documentclass[11pt]{article}
\usepackage{373,handout,url,graphicx}
%\usepackage{myconcrete,euler}
%\hidesolutions

% ======================================================================
\begin{document}

%\thispagestyle{empty}

\headers{CS 373}{Homework 4}{Spring 2003}
\begin{center}
{\LARGE\bf      CS 373:  Algorithms, Spring 2003}
\\[1ex]
{\Large\bf      Homework 4}
\\[0.25in]
\LARGE
\end{center}

%----------------------------------------------------------------------

\fbox{
\begin{minipage}{0.9\linewidth}
    ``At the sight of the still intact city, he remembered
    his great international precursors and set the whole
    place on fire with his artillery in order that those who
    came after him might work off their excess energies in
    rebuilding.''
    \\
    \hspace*{2cm} -- The tin drum, Gunter Grass
\end{minipage}
}


\section{Required Problems}


\begin{enumerate}    
    \item \textsc{Some number theory.}\\
    \points{10} 
    \begin{enumerate}
        \item \points{5} Prove that if $\gcd(m,n)=1$, then
        $m^{\phi(n)} + n^{\phi(m)} \equiv 1 (\bmod m n)$.
        \begin{proof}
            From $gcd(m,n) = 1$ and Euler's Theorem,
            \begin{eqnarray*}
                m^{\phi(n)} &\equiv& 1\ (mod\ n)\\
                n^{\phi(m)} &\equiv& 1\ (mod\ m)
            \end{eqnarray*}
            holds, and $m^{\phi(n)}$ and $n^{\phi(m)}$ are
            represented as $(0, 1)$ and $(1, 0)$ in
            $\Z_m\times \Z_n$, respectively. Therefore, we
            have
            \begin{eqnarray*}
                m^{\phi(n)} + n^{\phi(m)} = (0, 1) + (1, 0) = (1, 1)
            \end{eqnarray*}
            in $\Z_m\times \Z_n$. On the other hand, from Chinese
            Remainder Theorem, 1 in $\Z_{mn}$ corresponds to $(1, 1)$ in
            $\Z_m\times\Z_n$. Thus 
            \begin{eqnarray*}
                m^{\phi(n)} + n^{\phi(m)} \equiv 1\ (mod\ m n)
            \end{eqnarray*}
        \end{proof}
        \item \points{5} Give two distinct proofs that there
        are an infinite number of prime numbers.
        \begin{itemize}
            \item
            \begin{proof}               
                Assuming there are finite number of
                prime numbers, and let us denote the
                $i_{th}$ prime number as $p_i$. Let
                the (product of all the primes - 1)
                be $P = \Prod{p_i} + 1$. Clearly, $P
                \mod p_i \ne 0$. Namely, $P$ is not
                dividable by any prime number.
                Namely, it is prime. A
                contradiction.
            \end{proof}  
       
            \item
            \begin{proof}
                First we will show that the Fermat numbers $F_n =
                2^{2^n} + 1$ are pairwise relatively prime. From
                induction, it is easy to prove that
                \begin{eqnarray*}
                    F_m - 2 = F_0 \times F_1 \cdots F_{m-1}.
                \end{eqnarray*}
                This implies that if $d$ divides both $F_n$ and
                $F_m$, $(n < m)$, then d also divides $F_m - 2$,
                that is, d divides 2. But every Fermat number is
                odd, thus d has to be 1. Therefore, any Fermat
                numbers are relatively prime.
                
                Here, choose a prime divisor $p_n$ of each Fermat
                number $F_n$. However, from the property of Fermat
                numbers given above, those primes are all
                distinct. Therefore there are infinite number of
                primes.
            \end{proof}
        \end{itemize}
    \end{enumerate}
    
    \item \textsc{Even More Number Theory}\\
    \points{10}
    
    Prove that $\cardin{P(n)} =\Omega(n^2)$, where $P(n) =
    \brc{(a,b)\sep{ a,b \in \ZZ}, 0<a < b \leq n,
       \gcd(a,b)=1}$.
    
    \begin{proof}
        Let $Q(n, b) = \brc{(a, b) \sep{ a \in \Z, 0 < a < b
              \le n, \gcd(a, b) = 1}}$. Clearly,
        $\cardin{Q(n,b)}=\phi(b)$, and $|P(n)| =
        \Sum_{i=2}^n Q(n, i)= \Sum_{i=2}^n \phi(i)$.
        Namely, we want to prove that $\Sum_{i=2}^{n}
        \phi(i) = \Omega(n^2)$.
        
        \begin{figure}[t]
            \begin{center}
                \includegraphics[width=5cm]{figs/grid}
            \end{center}
            \caption{Grids and a line from the origin.}
            \figlab{grid}
        \end{figure}

        
        Consider the following set of pairs $R(n) =
        \brc{(a,b)|a,b\in \Z, 0<a<b\le n}$.  Clearly,
        $\cardin{R(n)} = \frac{n(n-1)}2$. Pairs in $R(n)$
        can be represented as points in 2D space, and we can
        draw straight lines from an origin $(0,0)$ to every
        point in $R(n)$. One observation is that, starting
        from an origin, a node $(a, b)$ that every straight
        line meets for the first time is in $P(n)$. It is
        simply because all the points on one straight line
        is written down as $(a\cdot c, b\cdot c)$, while $a$
        and $b$ are relatively prime and $c$ is a positive
        integer.
        
        Therefore, total number of points on this straight
        line, defined by a pair $(a,b)$, is 
        $\floor{n/b}$. Thus,
        \begin{eqnarray*}
            \Sum_{b=2}^n \pth{\floor{n/b}\cdot |Q(n, b)| \MakeBig} 
            = \Sum_{b=2}^n \pth{\floor{n/b}\cdot
               \phi(b)\MakeBig} 
            = \frac{n(n-1)}2.
        \end{eqnarray*}
        We overestimate $\floor{n/b}$ as $n/b$, and remove
        $n$ from both sides, then we have $\frac{n(n-1)}2
        \le \Sum_{b=2}^n \pth{n/b\cdot \phi(b)}$. Namely,
        \begin{eqnarray*}
            \frac{n-1}2 &\le& \Sum_{b=2}^n \frac{\phi(b)}b
            = \Sum_{b=2}^{\floor{n/4}} \frac{\phi(b)}b 
            + \Sum_{b=\floor{n/4}+1}^n \frac{\phi(b)}b
            \le \frac{n}4 + \Sum_{b=\floor{n/4}+1}^n
            \frac{\phi(b)}b\\
            &\le& \frac{n}4 + \Sum_{b=\floor{n/4}+1}^n \frac{\phi(b)}{n/4}
            \le \frac{n}4 + \Sum_{b=2}^n \frac{\phi(b)}{n/4},
        \end{eqnarray*}
        since $\phi(n) < n)$.  Thus, $\frac{n-2}4 \le
        \Sum_{b=2}^n \frac{\phi(b)}{n/4}$, which implies
        that
        $\frac{n(n-2)}{16} \le \Sum_{b=2}^n \phi(b) = |P(n)|$.
        Namely, $|P(n)| = \Omega(n^2)$.
    \end{proof}
    
    
    \item \textsc{Yet Another Number Theory Question}\\
    \points{20}
    
    \begin{enumerate}
        \item \points{2} Prove that the product of all
        primes $p$, for $m <p \leq 2m$ is at most
        $\binom{2m}{m}$.
        
        \begin{proof}
            Let us denote all the primes between $m$ and $2m$ as $p_i$, and
            all the other numbers as $n_i$. Then the following holds
            \begin{eqnarray*}
                {2m\choose{m}} &=& \frac{2m\cdot(2m -
                   1)\cdots(m + 2)
                   \cdot(m + 1)}{m\cdot(m - 1)\cdots{2}\cdot{1}}
                = \frac{\Prod{p_i}\cdot\Prod{n_i}}{m\cdot(m
                - 1)\cdots{2}
                \cdot{1}}.
            \end{eqnarray*}
            Since none of the numbers between $2$ and $m$
            divides any of the prime $p_i$, it must be that
            the number $\frac{\Prod{n_i}}{m\cdot(m -
               1)\cdots{2}\cdot{1}}$ is an integer number,
            as ${2m\choose{m}}$ is an integer. Therefore, by
            using a positive integer $c$, we have
            ${2m\choose{m}} = c\cdot \Prod{p_i}$   , that is,
            product of all the primes between $m+1$ and $2m$
            is at most ${2m}\choose{m}$.
        \end{proof}
        
        \item \points{4} \textbf{Using (a)}, prove that the
        number of all primes between $m$ and $2m$ is
        $O(m/\ln m)$.
        \begin{proof}
            Let us denote all primes between $m$ and $2m$ as
            $p_1 < p_2 < \cdots< p_k$.  Since $p_1 \geq m$,
            it follows by (a) that $m^k \le \Prod_{i=1}^{k}
            {p_i} \le {2m\choose{m}}\leq 2^{2m}$.
            Now, taking log of both sides, we have
            $k \lg{m} \leq 2m$. Namely, $k \leq 2m/\lg m$.
        \end{proof}

        \item \points{3} \textbf{Using (b)}, prove that the
        number of primes smaller than $n$ is $O(n/\ln n)$.
        \begin{proof}
            Let the number of primes less than n be
            $\Pi(n)$, then from (b) there exist some
            positive constant $C$, such that for all
            $\forall n \ge N$, we have $\Pi(2n) - \Pi(n) \ge
            C\cdot n/\ln{n}$. Namely, $\Pi(2n) \geq C\cdot
            n/\ln{n}+ \Pi(n)$. Thus,
            \[
            \Pi(2n) \geq \sum_{i=0}^{\ceil{\lg n} } C\cdot
            \frac{n/2^i}{\ln \pth{n/2^i}} = \Omega\pth{\frac{n}{\ln{n}}},
            \]
            by observing that the summation behaves like a
            decreasing geometric series (verify this!).
        \end{proof}
        
        \item \points{2} Prove that if $2^k$ divides
        $\binom{2m}{m}$ then $2^k \leq 2m$.

        \noindent{\bf Solution}:
            Refer to (e)

        \item \points{5} (Hard) Prove that for a prime $p$, if $p^k$
        divides $\binom{2m}{m}$ then $p^k \leq 2m$.
        
        \begin{proof}
            The proof is by induction.  Base case: When $m =
            1$, we have
            \begin{eqnarray*}
                p^k | {2m\choose{m}} &\rightarrow& p^k | {2\choose{1}} = 2\\
                &\rightarrow& p = 2 \ and\ k = 1\\
                &\rightarrow& p^k = 2 \le 2m = 2
            \end{eqnarray*}
            
            {\bf Inductive step:} Any $m$ can be described
            as $m=p\beta+\alpha\ (0 \le \alpha < p)$.
            Consider a situation where $p^k|{2m\choose{m}}$.
            \begin{eqnarray*}
                {2m\choose{m}}&=&\frac{2m(2m-1)\cdots 
                   (m+2)(m+1)}{m(m-1)\cdots 2\cdot 1}\\
                &=&\frac{(2p\beta + 2\alpha)(2\beta +
                   2\alpha - 1)
                   \cdots (p\beta+\alpha +2)(p\beta + \alpha+1)}
                {(p\beta+\alpha)(p\beta+\alpha-1)\cdots 2\cdot 1}.
            \end{eqnarray*}
            Since only a factor of $p$ matters in this division, we rewrite
            $2m\choose{m}$ by leaving only multiples of
            $p$. Formally, let 
            \[
            B(m,p) =
            \pth{\Prod_{i=m+1, p | i}^{2m} i}
            / \pth{\Prod_{i=1, p | i}^{m} i}.
            \]
            Assume for the time being that $B(m,p)$ is an
            integer number, and observe that if $p^k |
            \binom{2m}{m}$ then it divides also $B(m,p)$,
            and $B(m,p) | \binom{2m}{m}$.  

            So, what is $B(m,p)$? Well, first consider
            the case when $\alpha < p/2$:
            \begin{eqnarray*}
                B(m,p) = \frac{(p\cdot 2\beta) \cdot (p\cdot
                   (2\beta-1)) \cdots (p\cdot(\beta + 1))}
                {(p\cdot\beta)\cdot (p\cdot(\beta-1))
                   \cdots (p\cdot 2) \cdot (p \cdot 1)}
                =\binom{2\beta}{\beta},
            \end{eqnarray*}
            which proves that $B(m,p)$ is an integer number
            in this case, and also that $p^k |
            \binom{2\beta}{\beta}$. By induction, since
            $\beta < m$, we have that $p^k \leq 2\beta \leq
            2m$, which establish the claim.
            
            For $\alpha \geq p/2$, we have
            \begin{eqnarray*}
                B(m,p) = 
                \frac{p\cdot (2\beta + 1)\cdot
                   p\cdot 2\beta\cdots p\cdot(\beta
                   + 1)}
                {p\cdot\beta\cdot
                   p\cdot(\beta-1)\cdots p\cdot 2
                   \cdot p} = p(2\beta + 1 ) \binom{2\beta}{\beta},
            \end{eqnarray*}
            which proves that $B(m,p)$ is an integer
            number. Now, if $p$ does not divide $2\beta +
            1$, then $p^{k-1}$ must divide $B(m,p)$, as such
            $p^{k-1} \leq 2\beta$. And $p^k \leq 2 \beta p
            \leq 2m$, establishing the claim.
            
            So, we are left with the case that $p$ divides
            $2\beta+1$ (which implies that $p>2$). But then,
            $p$ can not possibly divide $\beta+1$ and
            $2(\beta+1)$, and furthermore, we know that $p^k
            | B(m,p)$. Thus,
            \[
            \beta = 2B(m,p) = \frac{2p(\beta+1)}{p(\beta+1)}
            B(m,p) = p(\beta+1) \binom{2(\beta+1)}{\beta+1}.
            \]
            Now, we know that $p^k | \beta$, which implies
            by induction that $p^k \leq 2 (\beta+1)p$. Which
            is not good enough, as this is larger than the
            bound we require.
            
            Well, just a little bit more patience. If $k=1$,
            the claim trivially hold (as one can easily
            verify). So, we can assume that $k >1$. But then
            we are dealing with the case that $p |
            2\beta+1$. Thus, $p^k \leq (2\beta+1)p + p$.
            However, $p^2$ divides the left side of this
            inequality, but not the right side of this
            inequality ($p^2 | p^k$ since $k>1$ and $p^2 |
            (2\beta+1)p$ since $p | 2\beta+1$). It follows
            that $p^k < (2\beta+1)p + p$. Namely, $p^k \leq
            (2\beta +1)p \leq 2m$.
  
            And we are FINALLY done. (I am sure BTW that
            there is a much shorter solution. This is a
            solution I could come up with.)
        \end{proof}
        
        \item \points{4} \textbf{Using (e)}, prove that that
        the number of primes between $1$ and $n$ is
        $\Omega(n /\ln n)$.  (Hint: use the fact that
        $\binom{2m}{m} \geq 2^{2m}/(2m)$.)
        \begin{proof}
            Assume $2m\choose{m}$ have $k$ prime factors,
            and thus can be written down as
            \begin{eqnarray*}
                {2m\choose{m}} = \Prod_{i=1}^k{p_i^{n_i}},
            \end{eqnarray*}
            where $n_i$ is the number of factors of $p_i$ in
            $2m\choose{m}$. From (e), we know that $p_i^{n_i} \le
            2m$. Therefore, the following is derived:
            \begin{eqnarray*}
                \frac{2^{2m}}{2m} \le {2m\choose{m}} \le 
                \Prod_{i=1}^k{2m} = (2m)^k.
            \end{eqnarray*}
            By taking $\lg$ of both sides, we have
            \begin{eqnarray*}
                \frac{2m - \lg(2m)}{\lg(2m)} \leq k.
            \end{eqnarray*}
            Since $k$ is \emph{less} than the number of
            primes between $1$ and $2m$, the number of
            primes small than $2m+1$ is
            $\Omega(m/\ln{m})$.
        \end{proof}
    \end{enumerate}
    
    \item \textsc{And now for something completely
       different.}\\
    \points{10}
    
    Prove that the following problems are NPC or provide a
    polynomial time algorithm to solve them:
    \begin{enumerate}
        \item Given a directed graph $G$, and two vertices
        $u,v \in V(G)$, find the maximum number of edge
        disjoint paths between $u$ and $v$.
        
        {\bf Solution:}
        
        This problem has a polynomial time algorithm. The
        maximum number of edge disjoint paths between $u$
        and $v$ is exactly equal to the value of maximum
        flow, where every edge is assigned weight one. Thus,
        we can use Edmonds-Karp algorithm, which running
        time is $O(|V||E|^2)$.
        
        \begin{center}
            \begin{algorithm}
                \textsc{Edge Disjoint Path($G(V,E),u,v$)}\+\\
                Run Edmonds-Karp($G(V,E),u,v$)
            \end{algorithm}         
        \end{center}
        
        \item Given a directed graph $G$, and two vertices
        $u,v \in V(G)$, find the maximum number of
        \emph{vertex} disjoint paths between $u$ and $v$
        (the paths are disjoint in their vertices, except of
        course, for the vertices $u$ and $v$).
        
        {\bf Solution:}
        
        This problem has a polynomial time algorithm. We
        will convert an input graph at each vertex as is
        shown in a figure, again, assigning each edge
        capacity 1. Then in the newly created graph, no 2
        edge-disjoint paths can share an edge that
        corresponds to a vertex in the original graph. It is
        because every vertex (except sink and source) has
        either only one incoming edge or one out-going edge,
        which allows only one flow to pass that vertex.
        Therefore, edge-disjoint paths correspond to
        vertex-disjoint paths in the original graph,
        respectively, and vice versa.  Here we can use
        Edmonds-Karp algorithm again.
        
        To convert a graph, it takes $O(|E|+|V|)$ time,
        because we have to look at every vertex and every
        edge. Let us denote the number of edges and vertices
        of the new graph as $|E^\prime|$ and $|V^\prime|$,
        respectively. Then the followings hold, because we
        add $|V|-2$ new edges and vertices:
        \begin{eqnarray*}
            |E^\prime| &=& |E| + |V| - 2\\
            |V^\prime| &=& |V| + |V| - 2.
        \end{eqnarray*}
        Edmonds-Karp algorithm takes
        \begin{eqnarray*}
            O(|V^\prime||{E^\prime}|^2) = O((2|V| - 2)(|E|+|V|-2)^2) = O(|V|(|E|+|V|)^2)
        \end{eqnarray*}
        
        \begin{center}
            \begin{algorithm}
                \textsc{Vertex Disjoint Path (G(V,E), u, v)}\+\\
                For all $v\in V\setminus \brc{u,v}$\+\\
                Add a new edge as in the figure\-\\
                Let us call this new graph as $G^\prime(V^\prime, E^\prime)$\\
                Run Edmonds-Karp($G^\prime(V^\prime,E^\prime),u,v$)
            \end{algorithm}
        \end{center}             
        \begin{figure}[htbp]
            \begin{center}
                \includegraphics{figs/vertex-disjoint}
                \caption{How to add a edge to a vertex}
            \end{center}
        \end{figure}            
    \end{enumerate}
    
    \item \textsc{Minimum Cut}\\
    \points{10}
    
    Present a \emph{deterministic} algorithm, such that
    given an undirected graph $G$, it computes the minimum
    cut in $G$. How fast is your algorithm? How does your
    algorithm compares with the randomized algorithm shown
    in class?
    
    {\bf Solution:}
    
    The deterministic algorithm consists of picking any pair
    of vertices $(s,t)$, and computing the minimum cut that
    have $s$ on one side, and $t$ on the other. Here, we
    take every edge and replicate it twice, in both
    directions, and assign every edge weight one. We next
    compute the maximum flow between $s$ and $t$. Clearly,
    by the min-cut max flow, one of the cut we would compute
    is the minimum cut in the graph. 

    The running time of this algorithm is not exciting. In
    the worst case, it would take $O(m^2n \cdot n^2) =
    O(n^8)$ time. This is much worse than the randomized
    algorithm, that takes roughly $O(n^2 \log^3 n )$ time.
    

    \item \textsc{Independence Matrix}\\
    \points{10}
    
    Consider a $0-1$ matrix $H$ with $n_1$ rows and $n_2$
    columns. We refer to a row or a column of the matrix $H$
    as a line. We say that a set of $1$'s in the matrix $H$
    is \emph{independent} if no two of them appear in the
    same line. We also say that a set of lines in the matrix
    is a \emph{cover} of $H$ if they include (i.e.,
    ``cover'') all the $1$'s in the matrix. Using the
    max-flow min-cut theorem on an appropriately defined
    network, show that the maximum number of independent
    $1$'s equals the minimum number of lines in the cover.
    
    {\bf Solution:}
    
    Given a $n_1\times n_2$ matrix $M$, we build a directed
    bipartite graph as follows.
    \begin{itemize}
        \item [1] Make bipartite graph with $n_1$ nodes on
        the left $\brc{l_i}$ and $n_2$ nodes on the right
        $\brc{r_i}$.
        \item [2] Draw an edge from $l_i$ to $r_j$ if
        $M(i,j) = 1$
        \item [3] Put source node and draw edges to every
        $l_i$
        \item [4] Put sink node and draw edges from every
        $r_i$
    \end{itemize}
    From how we build a graph, we know that an edge in a
    bipartite graph corresponds to $1$ in a matrix $M$. When
    the given matrix is as follows, a bipartite graph looks
    like the one in the figure.
    \begin{eqnarray*}
        \left(
            \begin{tabular}[tb]{cccc}
                0 & 1 & 0 & 0 \\
                0 & 0 & 1 & 0 \\
                0 & 1 & 0 & 0 \\
                1 & 0 & 0 & 1 \\
            \end{tabular}
        \right)
    \end{eqnarray*}
    \begin{figure}[htbp]
        \begin{center}
            \includegraphics{figs/bipartite}
            \caption{Generated bipartite graph.}
        \end{center}
    \end{figure}
    
    We will show that flow corresponds to independent 1's and cut
    corresponds to a cover, in the following.
    \begin{itemize}
        \item Since there is only one path from source to
        each $l_i$ and from each $r_j$ to sink, a node in
        the bipartite graph can be shared by only one path.
        This implies that all the $1$s that correspond to
        flow must be independent. On the other hand, given
        independent $1$s, we can actually make flows by
        using all the edges that correspond to the $1$s.
        Therefore, flow is equivalent with independent $1$s.
        \item From the way how we build a graph, $i_{th}$ row
        corresponds to $l_i$ and $j_{th}$ column corresponds to
        $r_j$. The following is a relation between a cover and
        a cut $C=(S,T)$:
        \begin{itemize}
            \item $l_i$ belongs to $S$ if and only if $l_i$
            is not in a cover. This implies $l_i$ belongs to
            $T$ if and only if $l_i$ is in a cover.
            \item $r_j$ belongs to $S$ if and only if $r_j$
            is in a cover. This implies $r_j$ belongs to $T$
            if and only if $r_j$ is not in a cover.
        \end{itemize}
        In fact, the size of a cover is equal to the size of
        a cut. It is because all the edges inside of the
        bipartite graph are either coming out of $T$, going
        into $S$, or both. That is, no edges in the
        bipartite graph are coming from $S$ and going into
        $T$. Therefore, the size of a cut is (the number of
        $l_i$ in $T$) $+$ (the number of $r_j$ in $S$),
        which is equal to the size of a cover.
    \end{itemize}
    From max-flow min-cut theorem and these equivalences, we know
    that the maximum number of independent $1$ is equal to the
    minimum number of lines in the cover.
    
\end{enumerate} 
\end{document}




                                %-----------------------------------------------------------------------
